\section{Presentación de la empresa}

\subsection{Identificación y propuesta de valor}

audIT es una empresa de ingeniería de software, arquitectura de datos y ciberseguridad industrial
constituida en 2022, con casa matriz y centro de ingeniería en Viña del Mar y un laboratorio de
dispositivos embarcados en Santiago. Su especialidad es la continuidad operacional en
organizaciones cuya producción ocurre lejos de cualquier instalación propia y con conectividad
intermitente.

\begin{tablaAudit}{Identificación de la empresa proponente. Formulario T-6}{empresa}{p{3.6cm}p{10.4cm}}
\textbf{Campo} & \textbf{Contenido} \\ \midrule
Razón social & audIT \\
Giro & Consultoría informática, plataformas de datos en nube, arquitectura híbrida, soluciones de telemetría de terreno y auditoría de ciberseguridad \\
Misión & Convertir operaciones críticas de campo en flujos de datos auditables, continuos y seguros \\
Visión & Ser el socio tecnológico de referencia en el Cono Sur para sistemas de misión crítica bajo arquitecturas resilientes que operan primero sin conexión \\
Valores & Rigor metodológico, seguridad por diseño, transparencia auditable y continuidad de servicio \\
Casa matriz & Viña del Mar, Región de Valparaíso \\
Laboratorio & Santiago, Región Metropolitana \\
\end{tablaAudit}

\subsection{Líneas de negocio}

La capacidad instalada se organiza en tres líneas que cubren el ciclo completo que este caso
requiere, desde el firmware del equipo en cabina hasta el tablero que usa la gerencia de finanzas.

\begin{enumerate}
    \item \textbf{Ingeniería de sistemas de terreno.} Desarrollo de firmware industrial,
    integración de telemetría vehicular multimarca sobre SAE J1939, rFMS y tacógrafos digitales, y
    provisión del paquete embarcado \emph{audIT EdgeHub}, que aporta almacenamiento local
    inalterable y sincronización automática tras desconexión prolongada.
    \item \textbf{Plataformas de datos y analítica.} Modelamiento e implantación de repositorios
    operacionales y analíticos en nube híbrida bajo el nombre \emph{audIT TelemetryCore}, ingestión
    masiva de eventos, motores de asignación y tableros de control operacional.
    \item \textbf{Arquitectura híbrida y ciberseguridad operacional.} Diseño de capas
    anticorrupción para convivencia con sistemas de gestión heredados, auditoría de seguridad
    perimetral y gobierno de datos personales bajo la Ley N.\textordmasculine{} 21.719
    \parencite{ley21719}.
\end{enumerate}

\subsection{Capacidad de intervención en terreno}

El Caso impone una condición que ninguna oferta puede resolver desde un escritorio. El equipamiento
a bordo solo puede instalarse, actualizarse o reemplazarse cuando el camión pasa por un terminal, y
los terminales del mandante se reparten a lo largo de más de 2.500 kilómetros entre Antofagasta y
Puerto Montt \parencite{bases_caso10_2026}. audIT mantiene convenios vigentes con una red técnica
certificada con presencia permanente en los cuatro nodos regionales del mandante, con capacidad de
intervención física en sitio y reemplazo de equipamiento telemático en menos de cuatro horas. Esta
capacidad da cumplimiento a lo exigido en RT-08.04 y RT-21.16 \parencite{bases_transversales_2026}.

\subsection{Acreditaciones pertinentes al caso}

RT-15.02 del Caso exige del adjudicatario conocimiento acreditado del régimen especial de jornada
del conductor de carga y del reglamento de transporte de sustancias peligrosas, además de
experiencia comprobable en telemática de flotas y en soluciones con operación desconectada
prolongada \parencite{bases_caso10_2026}. audIT declara esas competencias y las respalda con la
experiencia consignada en el Formulario T-6 que acompaña esta oferta.
